\documentclass{beamer}
\def\presentationtype{0}
\def\corsodilaurea{Ingegneria Informatica e dell'Automazione Industriale}
\def\insegnamento{Ottimizzazione}
\def\nolezione{8}
\def\titolo{Il metodo del simplesso -- Il ``tableau'' e l'operazione ``pivot''}
\def\noattivita{1}

\def\argomenti{Programmazione Lineare}
\def\obiettivilezione{Definire il ``tableau'' e l'operazione ``pivot'' per poter sviluppare metodo del simplesso primale}
\def\nucleotematico{Richiami}
\def\contenuto{Programmazione lineare}


\usepackage[italian]{babel}
%
% Hyperref info for PDF
%
\usepackage{hyperref}
\hypersetup{%
%pdfpagemode=FullScreen,
pdfpagemode=FullScreen,
pdfstartview=Fit,
pdfsubject={\argomenti},
pdfkeywords={\insegnamento,\titolo}
}

\definecolor{ecampus@engred}{RGB}{227,28,23}
\setbeamercolor{normal text}{fg=black,bg=}
\setbeamercolor{alerted text}{fg=red}
\setbeamercolor{example text}{fg=green!50!black}
\setbeamercolor{structure}{fg=ecampus@engred,bg=}
\usetheme{default}
\useinnertheme{rounded}%{circles}
\setbeamertemplate{blocks}[default]
\setbeamercolor{block title}{bg=}
\setbeamercolor{block body}{bg=}

\ifnum\presentationtype>0 {\setbeamertemplate{enumerate item}{Es. \nolezione.\presentationtype.\arabic{enumi})}} \fi

%\setbeamertemplate{frametitle continuation}[from second][(cont'd)]
\usefonttheme[stillsansseriftext,stillsansserifsmall]{serif}
\setbeamertemplate{headline}{%
\leavevmode%
   \hspace*{4.75cm}
   \scalebox{0.90}{%
    \begin{minipage}{\linewidth}%
      \tiny
 	\begin{tabular}{ll}
	  \\
% 	  \textcolor{gray}{Classe}					& \corsodilaurea\\
%	  \textcolor{gray}{Argomento:}					& \insegnamento\\
%	  \textcolor{gray}{Lezione n}:							& \nolezione \\
	  \textcolor{gray}{Titolo:}			& \titolo \\
	  \textcolor{gray}{Attivit\`a}:		& 	  \ifnum\presentationtype=0 {Lezione \nolezione} \else {Sessione di studio \nolezione.\presentationtype} \fi
	\end{tabular}%
    \end{minipage}%
   }%
}
\setbeamertemplate{footline}[text line]{%
  \begin{minipage}{\linewidth}%
    \hfill\insertpagenumber%
%    
    \begin{center}%
      \rule{\linewidth}{0.4pt}
    \end{center}
    \vspace*{-0.65cm}
    \begin{center}%
      \scalebox{0.55}{%
	\begin{minipage}{\linewidth}%
	  \begin{center}%
	    {\tiny
	    \copyright \ 2022 Gionata Massi - \color{blue}{\href{mailto:gionata.massi@savoiabenincasa.it}{gionata.massi@savoiabenincasa.it}}\\
	    \phantom{XXX}}
	  \end{center}
	\end{minipage}%
      }%
    \end{center}%
%
  \end{minipage}%
  }
\setbeamertemplate{navigation symbols}{}

%% Save up on ink for the 4-up handouts
\mode<handout>{%
  \pgfpagesuselayout{4 on 1}[a4paper, landscape, border shrink=10mm]
  \pgfpageslogicalpageoptions{1}{border code=\pgfstroke}
  \pgfpageslogicalpageoptions{2}{border code=\pgfstroke}
  \pgfpageslogicalpageoptions{3}{border code=\pgfstroke}
  \pgfpageslogicalpageoptions{4}{border code=\pgfstroke}
}

\mode<presentation>{\AtBeginSection{%
  \begin{frame}
    \frametitle{Piano della presentazione}
    \tableofcontents[currentsection]
  \end{frame}}
}

\usepackage{microtype}
\usepackage[utf8]{inputenc}
\usepackage[T1]{fontenc}
%\usepackage[osfss]{libertine}
\usepackage{palatino}
\usepackage[scaled=.77]{beramono}
\usepackage{booktabs}
% \usepackage{attachfile}

% Toglie l'enorme spazione prima di description
%\setbeamersize{description width=0pt}

\usepackage{tikz}
\usetikzlibrary{decorations,arrows,shapes,backgrounds,matrix,positioning}
\usepackage{multirow}
\usepackage{verbatim}
\usepackage{eurosym}
\usepackage{pgfpages}

\usepackage{zref-savepos}
\newcounter{restofframe}
\newsavebox{\restofframebox}
\newlength{\mylowermargin}
\setlength{\mylowermargin}{2pt}

\newenvironment{restofframe}{%
    \par%\centering
    \stepcounter{restofframe}%
    \zsavepos{restofframe-\arabic{restofframe}-begin}%
    \begin{lrbox}{\restofframebox}%
}{%
    \end{lrbox}%
    \setkeys{Gin}{keepaspectratio}%
    \raisebox{\dimexpr-\height+\ht\strutbox\relax}[0pt][0pt]{%
    \resizebox*{!}{\dimexpr\zposy{restofframe-\arabic{restofframe}-begin}sp-\zposy{restofframe-\arabic{restofframe}-end}sp-\mylowermargin\relax}%
        {\usebox{\restofframebox}}%
    }%
    \vskip0pt plus 1filll\relax
    \mbox{\zsavepos{restofframe-\arabic{restofframe}-end}}%
    \par
}

\title{\titolo}
\author{Gionata Massi}
\date{\today}

\setbeamercolor*{frametitle}{fg=black}
\setbeamercolor*{title}{fg=black}
\setbeamerfont{frametitle}{shape=\scshape,family=\rmfamily,size=\large,series=\bfseries}

\setbeamersize{}
\makeatletter
\newcommand\thefontsize[1]{{#1 The current font size is: \f@size pt\par}}
\makeatother

\tikzstyle{nicebox}=[draw=gray!100, fill=blue!10, very thick,
rounded corners, rectangle, inner sep=4pt, inner ysep=16pt]
\tikzstyle{niceboxtitle}=[draw=gray!100, fill=white, text=black,
rounded corners, very thick, rectangle]
\newcommand\nicebox[2]{
{\centering
\begin{tikzpicture}
\node [nicebox](box){
\begin{minipage}{0.95\textwidth}\centering
\begin{minipage}{0.95\textwidth}
#2
\end{minipage}\end{minipage}};
\node[niceboxtitle, right=10pt] at (box.north west)
{\small\textbf{#1}};
\end{tikzpicture}\par}
}

\tikzstyle{modelbox}=[draw=structure!100, fill=white!50, very thick,
rounded corners, rectangle, inner sep=4pt, inner ysep=16pt, text=blue!50!black]
\tikzstyle{modelboxtitle}=[draw=structure!100, fill=white!50, text=blue!50!black,
rounded corners, very thick, rectangle]
\newcommand\modelbox[2]{
{\centering
\begin{tikzpicture}
\node [modelbox](box){
\begin{minipage}{0.95\textwidth}\centering
\begin{minipage}{0.95\textwidth}
#2
\end{minipage}\end{minipage}};
\node[modelboxtitle, right=10pt] at (box.north west)
{\small\textbf{\mbox{#1}}};
\end{tikzpicture}\par}
}

% % %
% Tipi di sessioni di studio
% % %
\def\domande{Domande}
\def\approfondimenti{Approfondimenti}
\def\esercizi{Esercizi}
\def\esempi{Esempi svolti}
\def\relazioni{Relazioni}
\def\attivitapratica{Attivit\`a pratica}
% % %

\newcommand{\generatitolo}{ %
\begin{frame}[plain]
\pdfbookmark{\ifnum\presentationtype=0 {Lezione \nolezione} \else {Sessione di studio \nolezione.\presentationtype} \fi}{\ifnum\presentationtype=0 {Lezione \nolezione} \else {Sessione di studio \nolezione.\presentationtype} \fi}
\noindent\hspace*{3.715cm}
   \scalebox{0.90}{%
    \begin{minipage}{\linewidth}%
      \tiny
 	\begin{tabular}{ll}
	  \\
% 	  \textcolor{gray}{Classe}					& \corsodilaurea\\
%	  \textcolor{gray}{Argomento:}					& \insegnamento\\
%	  \textcolor{gray}{Lezione n}:							& \nolezione \\
	  \textcolor{gray}{Titolo:}			& \titolo \\
	  \textcolor{gray}{Attivit\`a}:		& 	  \ifnum\presentationtype=0 {Lezione \nolezione} \else {Sessione di studio \nolezione.\presentationtype} \fi
	\end{tabular}%
    \end{minipage}%
   }%
   
  \begin{center}
    \large{\textbf{\insegnamento}}
  \end{center}

  \vspace*{0.5cm}
  
  \begin{center}
    \huge{\titolo}
  \end{center}
  
  \begin{center}
	  \ifnum\presentationtype=0
	    {Lezione n$^{\circ}$ \nolezione}
	  \else
	    {Sessione di studio n$^{\circ}$ \nolezione .\presentationtype}
	  \fi
  \end{center}
  
  \vspace*{0.5cm}

  \begin{center}
    \large{\textsl{Gionata Massi}}\\
    \scriptsize\color{blue}{<\href{mailto:gionata.massi@savoiabenincasa.it}{gionata.massi@savoiabenincasa.it}>}
  \end{center}
  
  \newcounter{pageleft}
  \zsaveposy{pageleft}
  \vspace*{\dimexpr\zposy{pageleft}sp-1.67cm}

  \begin{minipage}{\linewidth}%
   \hfill\phantom{\insertpagenumber}%

  \begin{center}%
      \rule{\linewidth}{0.4pt}
    \end{center}
    
    \vspace*{-0.95cm}
    
    \begin{center}%
      \scalebox{0.55}{%
	\begin{minipage}{\linewidth}%
	  %\begin{center}%
	  %  \tiny
	  %  \copyright \ 2022 Gionata Massi - \color{blue}{\href{mailto:gionata.massi@savoiabenincasa.it}{gionata.massi@savoiabenincasa.it}}\\
	  %  \color{blue}{\href{mailto:gionata.massi@savoiabenincasa.it}{gionata.massi@savoiabenincasa.it}}\\
	  %  \phantom{XXX}
	  %\end{center}
	\end{minipage}%
      }%
    \end{center}%
     \phantom{XXX}
%
  \end{minipage}%

\end{frame}
}

\renewcommand{\vec}{\mathbf}
\newcommand{\matr}{\mathbf}

\uselanguage{italian}
\languagepath{italian}
\deftranslation[to=italian]{Theorem}{Teorema}
\deftranslation[to=italian]{theorem}{teorema}
\deftranslation[to=italian]{Definition}{Definizione}
\deftranslation[to=italian]{definition}{definizione}
\deftranslation[to=italian]{Corollary}{Corollario}
\deftranslation[to=italian]{corollary}{corollario}

\let\definition\relax
\let\enddefinition\relax

%\theoremstyle{example}
\newtheorem{definition}[theorem]{\translate{Definition}}

\def\fieldN{\mathbb{N}}
\def\fieldR{\mathbb{R}}
\def\fieldZ{\mathbb{Z}}

\def\matrA{\matr{A}}
\def\vecB{\vec{b}}
\def\vecC{\vec{c}}
\def\vecX{\vec{x}}

\usepackage{cancel}
\usepackage{bookmark}

% tkiz ball item
\newcommand*\circled[1]{\raisebox{2.5pt}{\tikz[baseline=(char.base)]{
            \node[circle,ball color=purple, shade, 
 color=white,inner sep=1.2pt] (char) {\tiny #1};}}}

% tkiz rounded item
\newcommand*\rounded[1]{\tikz[baseline=(char.base)]{
            \node[draw=none,ball color=purple, shade, 
 color=white, rounded corners=3.5pt, inner sep=2.5pt] (char) {\scriptsize #1};}}

 \usepackage{etoolbox}
\makeatletter
\apptocmd{\beamer@@frametitle}{\write\@auxout{\string\@writefile{frm}{\string\frametitleentry{\the\c@framenumber}{#1}{#2}}}}{}{}
\newcommand*{\frametitleentry}[3]{\@namedef{frametitleshort#1}{#2}\@namedef{frametitle#1}{#3}}
\AtEndDocument{\if@filesw\newwrite\tf@frm\immediate\openout\tf@frm\jobname.frm\relax\fi}
\@input{\jobname.frm}
\makeatother
 
%\usebackgroundtemplate{\includegraphics[width=\paperwidth]{../img/sfondo_ecampus.png}}

\usepackage{listings}
% Extension for listings package
\lstdefinelanguage{ampl}{
  % See AMPL book (1993 edition), Appendix A
  morekeywords={%
    % Table A.1
    % Excluding symbols and "not X", where X (and not) are keywords
    if, then, else, or,  exists, forall,  and, in, within, 
    not, union, diff, symdiff, inter, cross, setof, by, 
    % Skipping: less,  sum, prod, min, max, div, mod,
    % Skipping built-in functions (Tables A-2 and A-3)
    % Section A.5 -- Declarations of model entities (subject to ->
    % subject, to)
    set, param, var, arc, minimize, maximize, subject, to, node,
    % Section A.6 -- Set declarations (not repeating set or within)
    dimen, default, card, ordered, by, reversed, circular,
    % Section A.7 -- Parameter declarations (not repeating in, default)
    param, binary, integer, symbolic, check, Infinity
    % Section A.8 -- Variable declarations
    var,  coeff, cover, obj,
    % Section A.13 -- Command language
    % Skipping most commands!
    include, model, data, let, objective, drop, restore, solve,
    solution, quit, end, option, reset, update, commands,
    % A.14 -- Imported functions
    function, pipe, 
    % A.17 synonyms for previously-defined keywords
    coef, difference, dimen, dimension, intersect, intersection, s.t., 
    % 
    % must be defined somewhere ...
    for, while, repeat, break,
    % NEW language features (http://www.ampl.com/NEW/newlang.html)
    % Grouped by page
    purge, problem, redeclare, xref,
    complements, 
    table, IN, OUT, INOUT, read, write, 
    expand,
    for, commands, repeat, until, while, break, continue,
  },
  sensitive=true,
  comment=[l]{\#},
  morecomment=[s]{/*}{*/},
  string=[b]",
  morestring=[b]',
}[keywords,comments,strings]

\definecolor{comment_c}{RGB}{60,128,49}
\definecolor{rowname_c}{RGB}{128,60,49}
\definecolor{reserved_c}{RGB}{49,60,128}

\lstset{
	language=ampl,%
	commentstyle=\color{comment_c},%
	backgroundcolor=\color{blue!15},%
    basicstyle=\small\ttfamily,%
    %emph={s.t.},emphstyle={\bfseries},%
    otherkeywords={s.t.},%
    numbers=left, numberstyle=\tiny, stepnumber=1, numbersep=5pt,%
}

\definecolor{col_set}{RGB}{128,0,0}
\definecolor{col_par}{RGB}{91,0,91}
\definecolor{col_idx}{RGB}{0,91,91}
\definecolor{col_var}{RGB}{0,0,128}

\def\insieme{%1}{\rm \textcolor{col_set}{%1}}}


\begin{document}

\generatitolo

\section{Forma standard della Programmazione Lineare}

\begin{frame}[allowframebreaks]{Forma standard della PL}

{\scriptsize Ogni problema di Programmazione Lineare pu\`o essere espresso in una forma, detta \alert{forma standard}:

\vspace*{0.4cm}

\modelbox{Forma standard della PL} %
{\[
\begin{array}{crrrcccrcr}
\min z=& c_1 x_1 &+& c_2 x_2    &+&\cdots&+& c_n x_n &+& d\\
{\rm s.t.} & a_{11} x_1 &+& a_{12} x_2 &+&\cdots&+& a_{1n} x_n & = & b_1\\
                & a_{21} x_1 &+& a_{22} x_2 &+&\cdots&+& a_{2n} x_n & = & b_2\\
           &          & &           &  &\ddots&   &                   \\
           &  a_{m1} x_1 &+&  a_{m2} x_2 &+&\cdots&+& a_{mn} x_n & = & b_m\\
           \multicolumn{10}{c}{x_1 \geq 0,\cdots\ x_n \geq 0} 
\end{array}
\]}
}

\framebreak

\begin{itemize}
\item Il problema \`e di minimizzazione 
\item I vincoli sono tutti di uguaglianza 
\item Le variabili sono tutte non negative 
\end{itemize}


\modelbox{Forma standard della PL} %
{\[
\begin{array}{crrr}
\min z=& \vec{c}^T \vec{x} &+& d \\
{\rm s.t.} & \matr{A} \vec{x} & =& \vec{b} \\
&\multicolumn{3}{c}{\vec{x} \geq \vec{0}} 
\end{array}
\]}

\framebreak

\begin{itemize}
\item    $\vec{c}$ \`e un vettore di dimensione $n$ di parametri reali noti;
\item     $\vec{x}$ \`e un vettore di dimensione $n$ di variabili reali non negative (incognite);
\item     $\vec{b}$ \`e un vettore di dimensione $m$ dei parametri reali noti;
\item     $\matr{A}$ una matrice di dimensione $m$x$n$, con parametri reali noti e $m < n$ ($\infty$ soluzioni);
\item     $d$ \`e un parametro reale noto.
\end{itemize}

\framebreak

\begin{description}
\item [$c_j$] coefficienti dei costi variabili $\quad j=1,2,\ldots,n$

\item [$b_i$] coefficienti delle risorse o richieste $\quad  i=1,2,\ldots,m$

\item [$d$] costante di costo fisso

\item [$z$] funzione obiettivo che definisce il costo totale se il problema di partenza \`e di minimo, il profitto totale se il problema di partenza \`e di massimo

\item [$a_{ij}$] coefficienti tecnologici $\quad i=1,2,\ldots,m,\, j=1,2,\ldots,n$

\item [$x_{j}$] variabili decisionali $\quad j=1,2,\ldots,n$
\end{description}

\framebreak

\begin{description}
\item [Problema di minimo]
            un problema di P.L. in forma standard \`e espresso in forma di minimo

\item [Vincoli di eguaglianza]
		le equazioni dei vincoli devono essere espresse mediante relazioni di uguaglianza

\item [Vincoli di non negativit\`a delle variabili]
	le soluzioni ammissibili devono avere tutte le variabili non negative

\framebreak 

\item [Sistema consistente e non ridondante]
	\item  Si assume che il problema ammetta soluzioni. Tale condizioni sono verificate analiticamente se:
		\begin{itemize}
		\item                 $\rm rango ([A|b] ) = rango (A ) = m$, il che implica che esista almeno una soluzione e che il sistema non \`e ridondante, ossia non esistono vincoli combinazione lineare degli altri
		\item $m < n$, il che implica che esistono infinite soluzioni alternative
		\end{itemize}
\end{description}
\end{frame}

\section{Soluzioni di base}

\begin{frame}[allowframebreaks]{Soluzioni di base}

Sia dato un sistema di equazioni lineari $\matr{A}\vec{x}  = \vec{b}$ con $\matr{A} \in M_{m,n}(\mathbb{R}),\ m < n,\ \vec{b} \in \mathbb{R}^m$ e
$\rm rango ([A|b] ) = \rm rango (A ) = m$. Questo sistema verifica le ipotesi della forma standard della P.L.

\begin{definition}[soluzione]
  Ogni vettore $\vec{x}$ che soddisfa $\matr{A}\vec{x}  = \vec{b}$ si dice soluzione del sistema di equazioni lineari
\end{definition}

\begin{definition}[soluzione ammissibile]
    Una soluzione $\vec{x}$ si dice ammissibile (realizzabile) se soddisfa
    $\vec{x} \geq \vec{0}$
\end{definition}

\framebreak

\begin{definition}[regione ammissibile]
  L’insieme di tutte le soluzioni ammissibili si dice regione (insieme) ammissibile.
\end{definition}

\begin{definition}[soluzione di base]
  Si riordino le colonne di $\matr{A}$ in modo da partizionare la matrice in modo che $\matr{A} = [\matr{B|N}]$ dove $\matr{B}$ \`e una matrice invertibile di ordine $m$ e $\matr{N} \in M_{m,n-m}(\mathbb{R})$. Il sistema $\matr{A}\vec{x}  = \vec{b}$ \`e equivalente al sistema $\matr{B} \vec{x}_B + \matr{N} \vec{x}_N  = \vec{b}$.

 La soluzione $\vec{x}=\left[\begin{array}{c}\vec{x}_B\\ \vec{x}_N\end{array}\right]=\left[\begin{array}{c}\matr{B}^{-1}\vec{b}\\ \vec{0}\end{array}\right]$ si dice soluzione di base.
 \end{definition}

 \framebreak

\begin{definition}[soluzione di base ammissibile]
  Se una soluzione di base soddisfa $\vec{x}_B \geq \vec{0}$ \`e detta soluzione di base ammissibile.
 \end{definition}

\begin{definition}[matrice di base]
     La matrice quadrata $\matr{B} $ \`e detta matrice di base.
 \end{definition}

\begin{definition}[matrice non di base]
    La matrice $\matr{N}$ \`e detta matrice non di base.
 \end{definition}

\begin{definition}[variabili di base]
   Il vettore $\vec{x}_B$ \`e il vettore delle variabili di base.
\end{definition}
\framebreak
\begin{definition}[variabili non di base]
    Il vettore $\vec{x}_N$ \`e il vettore delle variabili non di base.
\end{definition}
 
 \begin{definition}[soluzione di base degeneri]
    Se una soluzione di base soddisfa $\vec{x}_B > 0$ allora $\vec{x}$ \`e detta soluzione non degenere. Se invece almeno una componente di $\vec{x}_B$ \`e nulla allora $\vec{x}$ \`e detta soluzione degenere. 
   \end{definition}
  
%  \begin{definition}[soluzione di base ammissibile]
{\small 
    Il numero di soluzioni di base per un problema di P.L. in forma standard con $n$ variabili e
   $m$ equazioni \`e al pi\`u uguale al numero di possibili scelte di $m$ su $n$ colonne, ossia il numero di
  soluzioni di base \`e limitato superiormente da
  \[ C(n ; m) = {n \choose m} = \frac{n!}{m! \cdot \left( n - m \right) !}\qquad n,k\in\mathbb{N}; 0\leq m \leq n \]
}

%   \end{definition}
  
\framebreak

  \begin{block}{Soluzioni di base - riepilogo}
  \begin{itemize}
  \item Una soluzione del sistema di equazioni lineari $\matr{A}\vec{x}  = \vec{b}$ ottenuta annullando un insieme di variabili non di base si dice soluzione di base.
  
   \item Una soluzione di base $\vec{x}$ che soddisfa $\vec{x} \geq \vec{0}$ \` e detta soluzione di base ammissibile.
  
   \item Una soluzione di base $\vec{x}$ in cui almeno una delle componenti di base \`e nulla \`e detta soluzione degenere. 
    
  \end{itemize}
   \end{block}
  
  \begin{definition}[soluzione ottima]
    Una soluzione di base ammissibile $\vec{x}^\star$ si dice ottima se verifica $\vec{c}^T\vec{x}^\star \leq  \vec{c}^T \vec{x}$ per ogni $\vec{x}$ soluzione ammissibile.
   \end{definition}
  
  \begin{definition}[soluzioni adiacenti]
    Due soluzioni di base $\vec{x} = \left[\begin{array}{c}\vec{x}_B\\ \vec{x}_N\end{array}\right]$ e $\vec{x}^\prime = \left[\begin{array}{c}\vec{x}_B^\prime\\ \vec{x}_N^\prime\end{array}\right]$ si dicono adiacenti se i vettori $\vec{x}$ e $\vec{x}^\prime$ contengono 
  le stesse variabili eccetto una.
 \end{definition}
\end{frame} 

\section{Teoremi fondamentali della Programmazione Lineare}

\begin{frame}[allowframebreaks]{Teoremi}

\begin{theorem}[esistenza delle soluzioni di base]
Se un sistema di equazioni lineari $\matrA\vecX  = \vecB$ ammette soluzioni ammissibili, ammette almeno una soluzione di base ammissibile
\end{theorem}

\begin{theorem}[equivalenza di punti estremi e soluzioni di base ammissibili]
   Sia $S$ il politopo convesso costituito da tutti i vettori $\vecX \in \fieldR^n$  che soddisfano $\matrA\vecX = \vecB, \vecX \geq \vec{0}$.
   
   Un vettore $\vecX$ \`e un punto estremo (vertice) di $S$ se e solo se $\vecX$ \`e una soluzione di base ammissibile di 
   $S := \{\matrA\vecX = \vecB, \vecX \geq \vec{0}\}$
\end{theorem}

\framebreak

\begin{block}{Corollario}
Se l’insieme convesso $S$ corrispondente al problema di P.L. in forma standard \`e non vuoto, allora esiste almeno un punto estremo
\end{block}

\begin{block}{Corollario}
L’insieme convesso $S$ corrispondente al problema di P.L. in forma standard possiede al pi\`u un numero finito di punti estremi
\end{block}

\begin{block}{Corollario}
 Se il politopo convesso $S$ corrispondente al problema di P.L. in forma standard \`e limitato allora $S$ \`e un poliedro convesso, ossia \`e costituito da punti che sono combinazione convessa di un numero finito di punti estremi
\end{block}

\framebreak

\begin{theorem}[Punti estemi e soluzione ottima]
  La funzione obiettivo lineare $z = \vecC^T\vecX + d$ con $\vecX$ appartenente al poliedro convesso $S$ raggiunge il minimo in corrispondenza di un punto estremo di $S$
\end{theorem}

\begin{theorem}[Soluzione ottima e soluzione di base]
    Se esiste una soluzione ammissibile ottima finita per il problema di P.L. in forma standard allora esiste almeno una soluzione di base ammissibile ottima
\end{theorem}

\end{frame}

\section{Forma canonica della Programmazione Lineare}

\begin{frame}{Forma canonica}
Un problema di P.L. in forma standard si dice in forma canonica rispetto ad una sequenza di base  $S = [s_1, s_2, \ldots, s_m]$ se:


\[\begin{array}{ll}
\alpha) & \matrA_S = [\matrA^{s_1} \matrA^{s_2} \cdots \matrA^{s_m}] = \matr{I}_m \\
\beta) & \vecC_S = \vec{0} \\
\gamma) & \vecB \geq \vec{0}
\end{array}
\begin{array}{ll}
\Bigg\} \text{ f.c. debole}\\
\phantom{x}
\end{array}
\left.\begin{array}{ll}
\phantom{x}\\
\phantom{x}\\
\phantom{x}\\
\end{array}
\right\rbrace \text{ f.c. forte}
\]

Se valgono $\alpha$) e $\beta$) il problema \`e in forma canonica debole. Se vale anche $\gamma$) il sistema \`e in forma canonica forte
\end{frame}


\begin{frame}[allowframebreaks]{Forma canonica: perch\'e?}

La condizione $\alpha$)  per cui $\matrA_S = [\matrA^{s_1} \matrA^{s_2} \cdots \matrA^{s_m}] = \matr{I}_m$  significa che il sistema $\matrA\vecX = \vecB$ \`e stato risolto rispetto alle variabili di base $x_{s_1}, x_{s_2}, \ldots, x_{s_m}$.
  

   La condizione $\beta$) per cui $\vecC_S = \vec{0}$ afferma che i coefficienti della funzione obiettivo $c_{s_1}, c_{s_2}, \ldots, c_{s_m}$ sono tutti nulli.

  La condizione $\gamma$) stabilisce che i termini noti sono tutti non negativi, ovvero che la soluzione di base \`e ammissibile.

\framebreak

{\footnotesize In corrispondenza di una forma canonica forte esiste una soluzione di base ammissibile data da: }

\small$$
\begin{array}{ll}
  x_{s_i} = b_i	& i=1,2,\ldots,m \\
  x_j = 0			&  j : j \notin S \\
  z = d				& \left( z= 
  \cancelto{0 \text{ per } c_{s_i }=0}{\displaystyle \sum_{i=1}^{m} c_{s_i} x_{s_i}} +
  \cancelto{0 \text{ per } x_j = 0}{\displaystyle \sum_{j:j\notin S} c_j x_j}
  +d \hspace*{0.8cm}
   \right)
\end{array}
$$
{\footnotesize A partire dalla forma canonica, attraverso l'operazione di pivot, sar\`a sviluppato il metodo del simplesso tableau}
\end{frame}

\end{document}
